%%%%%%%%%%%%%%%%%%%%%%%%%%%%%%%%%%%%%%%%%%%%%%%%%%%%%%%%%%%%%%%%%%%%%%%%%%%%%%%%%%
\begin{frame}[fragile]\frametitle{}
\begin{center}
{\Large निर्वाण षटकम्}
\end{center}

{\tiny (Ref: आदि शंकराचार्य --- निर्वाण षटकम् | Sanskrit Channel)}

\end{frame}

%%%%%%%%%%%%%%%%%%%%%%%%%%%%%%%%%%%%%%%%%%%%%%%%%%%%%%%%%%%
\begin{frame}[fragile]\frametitle{परिचय}

	\begin{itemize}
	\item आदि शंकराचार्यांनी रचलेले ६ श्लोकांचे स्तोत्र
	\item प्रत्येक श्लोकाचा अंत: \textbf{चिदानन्दरूपः शिवोऽहम् शिवोऽहम्}
	\item मुख्य तत्त्व: मी शरीर, मन, इंद्रिये, पंचमहाभूते किंवा पंचकोश नाही
	\item मी शुद्ध चैतन्य आणि आनंद आहे --- मी शाश्वत शिव आहे
	\end{itemize}

{\tiny (Ref: आदि शंकराचार्य --- निर्वाण षटकम्)}

\end{frame}

%%%%%%%%%%%%%%%%%%%%%%%%%%%%%%%%%%%%%%%%%%%%%%%%%%%%%%%%%%%
\begin{frame}[fragile]\frametitle{श्लोक १}

\begin{verse}
मनोबुद्ध्यहङ्कार चित्तानि नाहं \\
न च श्रोत्रजिह्वे न च घ्राणनेत्रे । \\
न च व्योम भूमिर्न तेजो न वायुः \\
चिदानन्दरूपः शिवोऽहम् शिवोऽहम् ॥१॥
\end{verse}

	\begin{itemize}
	\item मी मन, बुद्धी, अहंकार किंवा स्मृती नाही
	\item मी कान, त्वचा, नाक किंवा डोळे नाही
	\item मी आकाश, पृथ्वी, अग्नी, पाणी किंवा वारा नाही
	\item मी चैतन्य आणि आनंदाचे रूप आहे --- मी शाश्वत शिव आहे
	\end{itemize}

\end{frame}

%%%%%%%%%%%%%%%%%%%%%%%%%%%%%%%%%%%%%%%%%%%%%%%%%%%%%%%%%%%
\begin{frame}[fragile]\frametitle{श्लोक २}

\begin{verse}
न च प्राणसंज्ञो न वै पञ्चवायुः \\
न वा सप्तधातुः न वा पञ्चकोशः । \\
न वाक्पाणिपादं न चोपस्थपायु \\
चिदानन्दरूपः शिवोऽहम् शिवोऽहम् ॥२॥
\end{verse}

	\begin{itemize}
	\item मी श्वास किंवा पाच वायू नाही
	\item मी सप्तधातू किंवा पंचकोश नाही
	\item मी वाणी, हात, पाय, उपस्थ किंवा पायू नाही
	\item मी चैतन्य आणि आनंदाचे रूप आहे --- मी शाश्वत शिव आहे
	\end{itemize}

\end{frame}

%%%%%%%%%%%%%%%%%%%%%%%%%%%%%%%%%%%%%%%%%%%%%%%%%%%%%%%%%%%
\begin{frame}[fragile]\frametitle{श्लोक ३ व ४}

\begin{verse}
न मे द्वेषरागौ न मे लोभमोहौ \quad चिदानन्दरूपः शिवोऽहम् ॥३॥ \\
न पुण्यं न पापं न सौख्यं न दुःखं \quad चिदानन्दरूपः शिवोऽहम् ॥४॥
\end{verse}

	\begin{itemize}
	\item (श्लोक ३) माझ्यात द्वेष, राग, लोभ, मोह, मद, मत्सर नाही
	\item (श्लोक ३) मला धर्म, अर्थ, काम किंवा मोक्षाची इच्छा नाही
	\item (श्लोक ४) पुण्य, पाप, सुख, दुःख --- काहीच नाही
	\item (श्लोक ४) मला मंत्र, तीर्थ, वेद, यज्ञ यांची गरज नाही; मी भोक्ता नाही
	\item चिदानन्दरूपः शिवोऽहम् शिवोऽहम्
	\end{itemize}

\end{frame}

%%%%%%%%%%%%%%%%%%%%%%%%%%%%%%%%%%%%%%%%%%%%%%%%%%%%%%%%%%%
\begin{frame}[fragile]\frametitle{श्लोक ५}

\begin{verse}
न मृत्युर्न शङ्का न मे जातिभेदः \\
पिता नैव मे नैव माता न जन्मः । \\
न बन्धुर्न मित्रं गुरुर्नैव शिष्यं \\
चिदानन्दरूपः शिवोऽहम् शिवोऽहम् ॥५॥
\end{verse}

	\begin{itemize}
	\item मला मरणाचे भय नाही, जात-पात नाही
	\item मला वडील नाहीत, आई नाही --- मी कधीच जन्मलो नाही
	\item मी नातेवाईक, मित्र, गुरू किंवा शिष्य नाही
	\item मी चैतन्य आणि आनंदाचे रूप आहे --- मी शाश्वत शिव आहे
	\end{itemize}

\end{frame}

%%%%%%%%%%%%%%%%%%%%%%%%%%%%%%%%%%%%%%%%%%%%%%%%%%%%%%%%%%%
\begin{frame}[fragile]\frametitle{श्लोक ६}

\begin{verse}
अहं निर्विकल्पो निराकाररूपो \\
विभुत्वाच्च सर्वत्र सर्वेन्द्रियाणाम् । \\
न चासङ्गतं नैव मुक्तिर्न मेयः \\
चिदानन्दरूपः शिवोऽहम् शिवोऽहम् ॥६॥
\end{verse}

	\begin{itemize}
	\item मी द्वैतरहित आहे, माझे रूप निराकार आहे
	\item मी सर्वत्र अस्तित्वात आहे --- सर्व इंद्रियांना व्यापून आहे
	\item मी संलग्न नाही, बंदिवान नाही, मुक्त नाही --- मी परे आहे
	\item मी चैतन्य आणि आनंदाचे रूप आहे --- मी शाश्वत शिव आहे
	\end{itemize}

\end{frame}
